\documentclass[11pt]{article}
\usepackage[top=.79in,bottom=.79in,left=.86in,right=.86in]{geometry}
\usepackage{amsmath,amssymb,amsthm,booktabs,microtype,xcolor}
\usepackage[colorlinks=true,urlcolor=blue!45!black]{hyperref}
\newtheorem{theorem}{Theorem}
\newcommand{\F}{\mathbb F}
\newcommand{\PP}{\mathbb P}
\title{Seven-step graph towers:\\a strict bound at every prime}
\author{}
\date{27 November 2026}
\begin{document}
\maketitle
\vspace{-2em}
\begin{abstract}
We sharpen the universal spanning-tree valuation bound for seven-step
rank-two graph towers at every prime $p>14$. For $p\equiv1\bmod14$ we
give a simple graph tower with two directions of maximal depth.
Exact searches give the maxima at the five smallest eligible primes.
A separate argument sharpens the universal bound for every prime
number of steps and sufficiently large residue characteristic.
\end{abstract}

\section{The seven-step bound}
For $m$ integer steps $S=(s_i)_{i=0}^{m-1}$ spanning $\F_p^2$ modulo $p$,
the level-$n$ graph has vertices $(\mathbb Z/p^n\mathbb Z)^2$ and
one undirected edge from $v$ to $v+s_i$ for each pair $(v,i)$.
Let $\kappa_n$ be its spanning-tree count, and define
\[
 M_j(u)=\sum_i(u\cdot s_i)^{2j}\quad\text{over }\F_p,\qquad
 r(L)=\min\{j\ge1:M_j(L)\ne0\},\qquad A=2\sum_{L\in\PP^1(\F_p)}r(L).
\]
We use the \href{https://github.com/VladimirReshetnikov/ProveIt/blob/21c0252ffe8611baa0fbf528fbb003d6c42ca9af/Analysis/GeneralizedRiemannHypothesis/Research/grh_transfers_lubin_tate_and_graph_towers/graphs.tex}{ProveIt valuation identity}
$v_p(\kappa_n)=A(p^n-1)/(p-1)-2n$ for $p>2m$. Newton's identities
give $r(L)\le m$. If $r(L)=m$, the squared projections are the
roots of $T^m-c$ with $c\ne0$. These roots are distinct nonzero
squares, and dividing by any one root gives $2m\mid p-1$.
The \href{https://github.com/VladimirReshetnikov/ProveIt/blob/21c0252ffe8611baa0fbf528fbb003d6c42ca9af/Analysis/GeneralizedRiemannHypothesis/Research/grh_transfer_extensions_2026-11-26/sections/graph_extremum.tex}{first-moment lemma}
bounds the first nonzero moment polynomial index by
$k=\lfloor m/2\rfloor$. Zero residues can be discarded for the lemma.
If the first nonzero polynomial is $M_t$ and pointwise depth is at
most $d$, its at most $2t$ projective roots give
$A\le2t(p+1+2d-2t)$.

\begin{theorem}
For $m=7$ and every prime $p>14$, $A\le6(p+1)+38$,
improving the former $6(p+1)+48$. For $p\not\equiv1\bmod14$,
the bound improves to $6(p+1)+36$.
If $M_1=M_2=0$ and two directions have depth seven, those are the
only directions of depth above three, and $A=6(p+1)+16$.
This last case occurs in simple graphs for every $p\equiv1\bmod14$.
\end{theorem}

For $t\le2$, the bound above gives at most $4(p+1)+40$, which is
below $6(p+1)+36$. When $t=3$, there are at most six exceptional
directions. If none has depth seven, each adds at most six above
the depth-three baseline. If exactly one has depth seven, it adds
eight and the other five add at most six each. It remains to
handle two or more depth-seven directions with $t=3$.

\section{The seven-step certificate}
Take two depth-seven directions as coordinate axes. Each projection
row has squared entries in a coset of $\mu_7$. As squaring is
bijective on $\mu_7$, row scalings, column permutations and column
signs put the rows in the form $a_i=z^i$, $b_i=\epsilon_i z^{\sigma(i)}$,
where $z$ has order seven, $\sigma(0)=0$, and $\epsilon_0=1$.
The coefficient of $X^2Y^2$ in $M_2$ requires $f(z^2)=0$ for
$f(Z)=\sum_{i=0}^6 Z^{(i+\sigma(i))\bmod7}$.

\pagebreak[2]
The table gives the absolute norms of $f$ over $\mathbb Q(\zeta_7)$
for the $6!=720$ permutations. A norm can be computed as the
determinant of multiplication by $f$ on $\mathbb Z[Z]/(\Phi_7)$.
If $f_h$ is the coefficient of degree $h$, the entry $(i,j)$ for
$0\le i,j<6$ is $f_{(i-j)\bmod7}-f_{(6-j)\bmod7}$.
Evaluation at $z^2$ annihilates its image modulo $p$, so $p$
divides its determinant. The checker also computes the norm by
an exact product of the six conjugates.
\begin{center}
\small
\begin{tabular}{lrrrrrrrr}
\toprule
Absolute norm & 0 & 49 & 343 & 392 & 2744 & 3479 & 8281 & 117649\\
Multiplicity  &19 &441 &42 &126 &28 &42 &21 &1\\
\bottomrule
\end{tabular}
\end{center}

The only prime $p>14$ dividing a nonzero entry is $71$.
Elsewhere the norm is zero, so irreducibility of $\Phi_7$ forces
all seven coefficients of $f$ to be one. The cross moment
$\sum_i a_i b_i$ is now a signed sum of the seven distinct roots
of unity. Among sign polynomials with constant coefficient $+1$,
the absolute norms are $0,64,512$, with counts $1,49,14$. The only
case vanishing at an odd prime has every sign positive.
The $X^3Y$ and $XY^3$ coefficients in $M_2$ give the two additional
exponent lists $3i+\sigma(i)$ and $i+3\sigma(i)$ modulo seven.
The same norm table applies (relabel the indices or the outputs
of $\sigma$). Thus all three exponent lists must be permutations.
The 720-case test leaves just $\sigma(i)=ki$ for $k=1,3,5$; the
identity is not spanning. For $k=3,5$, geometric sums give
$M_3=105X^2Y^4$ or $105X^4Y^2$, respectively. As $p>14$ this
has only the two axis roots. Both have depth seven.

At $71$ take $z=20$ and check all $720\cdot64$ permutation/sign
choices by testing the mixed coefficients of $M_1$ and $M_2$.
After excluding dependent rows, precisely $k=3,5$ with positive
signs survive. This finishes the two-direction classification.
The steps $(z^i,z^{3i})$ in that classification span and have
distinct nonzero signed classes, so give simple graphs at every
level. Their depth pattern proves the claimed attainment.

\section{An obstruction for a prime number of steps}
\begin{theorem}
Let $m\ge5$ be prime, $k=(m-1)/2$, and
$p>\max(m^{m-1},27^k)$. Then
\[
 A\le 2k(p+1+2m-2k)-2(m-3).
\]
More precisely, if $M_1=0$ there are at most two depth-$m$
directions.
\end{theorem}
\begin{proof}
Suppose there are three distinct depth-$m$ directions. Choose a
primitive $m$th root $z\in\F_p$. As above, normalize the projection
rows to $a_i=z^i$, $b_i=\epsilon_i z^{\sigma(i)}$ and
$c_i=\eta_i z^{\tau(i)}$. Lift by replacing $z$ with a complex
primitive root $\zeta$. A cross moment $\sum_i a_i b_i$ and
all $3\times3$ minors of the three-row matrix have zero reduction
at $(p,\zeta-z)$. If one is nonzero, its integral norm is
divisible by $p$. Each conjugate of the cross moment has
absolute value at most $m$, and each conjugate of a minor at
most $3\sqrt3$ by Hadamard's inequality. The norm bounds are
therefore $m^{m-1}$ and $27^k$, both less than $p$. Thus the
cross moment and the minors vanish exactly.

Reducing its exponents modulo $m$, the cross moment is a sum of
$m$ signs on powers of $\zeta$. Its coefficient vector is
constant by irreducibility of $\Phi_m$. The coefficient sum is
odd with absolute value at most $m$, so the constant is $\pm1$
and all signs agree. Absorb the common sign in $b$.
The three rows are pairwise independent, as witnessed already
in their reductions modulo $p$. The vanishing minors give
$c=\alpha a+\beta b$ with $\alpha\beta\ne0$.
Every ratio $w=b_i/a_i$ satisfies $|w|=|\alpha+\beta w|=1$.
Expanding the second condition gives a real affine line in $w$;
it intersects the unit circle in at most two points.
Thus $\sigma$ acts by at most two translations $u,v$ modulo $m$.
If it translates $I$ by $u$ and the complement by $v$, bijectivity
gives $I+u=I+v$. A proper nonempty $I$ cannot be invariant
under a nonzero translation modulo a prime. Therefore only one
translation occurs, making $a,b$ proportional, a contradiction.

For $t=k$, there are at most $2k$ exceptional directions, and
all but two have depth at most $m-1$. Relative to the universal
bound the savings is at least $2(2k-2)=2(m-3)$. If $t<k$,
the universal bound for $t=k-1$ already saves $2(p+5)$ or more.
\end{proof}

\section{Five exact benchmarks}
To enumerate lists with $M_1=M_2=0$, normalize two independent
steps to the standard basis and use fixed sign representatives
for the other five. The three quadratic and five quartic
coefficients, after dividing by the binomial factors, are
linear in the sum of the keys
\[
  (x^2,xy,y^2,x^4,x^3y,x^2y^2,xy^3,y^4).
\]
Store every sorted pair by its key sum. For each sorted triple,
fetch every complementary pair whose first index is at least the
last triple index. This is an exhaustive search, including
repetitions and zero residues. At $29$ there are exactly twelve
normalized candidates. The best has $A=196$; lists with first
nonzero index at most two satisfy $A\le160$. One witness is the
basis and
\[
  (2,3),\ (6,2),\ (10,24),\ (11,21),\ (12,10).
\]
Its only exceptional directions are $(1,7)$ and $(1,24)$, both at
depth seven.

At $17,19,23,31$, a second pass uses just the first three key
coordinates, so exhausts every list with $M_1=0$. At these primes
the depth cap is six; lists with $M_1\ne0$ have $A\le2(p+1)+20$,
strictly below the best in each class. Here are the exact maxima:
\begin{center}
\begin{tabular}{@{}rrrr@{}}
\toprule
$p$ & Zero allowed & Nonzero residues & Simple graphs\\
\midrule
17 & 88  & 82  & 82\\
19 & 120 & 94  & 94\\
23 & 104 & 104 & 104\\
29 & 196 & 196 & 196\\
31 & 192 & 140 & 140\\
\bottomrule
\end{tabular}
\end{center}
The search logs give normal-form counts and explicit simple
witnesses. Here simplicity is enforced by distinct nonzero
residues up to sign, which ensures simple graphs at every layer.
The valuation identity turns each column into exact all-layer
maxima within the stated class.
For example, at $p=29$ the simple maximum is
$196(29^n-1)/28-2n$.

\paragraph{Scope.}
These conclusions concern the odd-step extremal problem in the
cited ProveIt reports. They are unconditional and do not
use RH or GRH. The exact seven-step maximum at general $p$ remains open.
\end{document}
